
\documentclass{article}
\usepackage{colortbl}
\usepackage{makecell}
\usepackage{multirow}
\usepackage{supertabular}

\begin{document}

\newcounter{utterance}

\centering \large Interaction Transcript for game `bbh', experiment `multistep\_arithmetic\_two', episode 93 with reflexion{-}Qwen3.5{-}4B.
\vspace{24pt}

{ \footnotesize  \setcounter{utterance}{1}
\setlength{\tabcolsep}{0pt}
\begin{supertabular}{c@{$\;$}|p{.15\linewidth}@{}p{.15\linewidth}p{.15\linewidth}p{.15\linewidth}p{.15\linewidth}p{.15\linewidth}}
    \# & \multicolumn{2}{c}{Player} && \multicolumn{2}{c}{Game Master} \\
    \hline

    \theutterance \stepcounter{utterance}  
    & & & \multicolumn{4}{p{0.6\linewidth}}{
        \cellcolor[rgb]{0.9,0.9,0.9}{
            \makecell[{{p{\linewidth}}}]{
                \texttt{\tiny{[P1$\langle$GM]}}
                \texttt{Solve multi{-}step arithmetic problems.} \\
\\ 
\texttt{Q: (({-}5 + 9 * {-}4 {-} 0) * (4 + {-}7 + 0 * {-}5)) =} \\
\texttt{A: 123} \\
\\ 
\texttt{Q: (({-}9 * 7 * 7 * {-}9) + (4 * {-}9 {-} 8 {-} {-}4)) =} \\
\texttt{A: 3929} \\
\\ 
\texttt{Q: (({-}3 + 5 * 8 * {-}4) {-} (9 {-} 8 * {-}7 + {-}9)) =} \\
\texttt{A: {-}219} \\
\\ 
\texttt{Q: (({-}9 + {-}3 + 6 {-} 8) {-} ({-}7 * {-}2 * {-}6 + {-}6)) =} \\
\texttt{A:} \\
            }
        }
    }
     \\ \\

    \theutterance \stepcounter{utterance}  
    & \multicolumn{4}{p{0.6\linewidth}}{
        \cellcolor[rgb]{0.561,0.737,0.561}{
            \makecell[{{p{\linewidth}}}]{
                \texttt{\tiny{[P1$\rangle$GM]}}
                \texttt{To solve the expression \$(({-}9 + {-}3 + 6 {-} 8) {-} ({-}7 * {-}2 * {-}6 + {-}6))\$, we will follow the order of operations (PEMDAS/BODMAS): Parentheses, Exponents, Multiplication/Division, Addition/Subtraction.} \\
\\ 
\texttt{**Step 1: Simplify the first part inside the outer parentheses.**} \\
\texttt{Expression: \$({-}9 + {-}3 + 6 {-} 8)\$} \\
\texttt{{-} Add \${-}9\$ and \${-}3\$: \${-}9 + ({-}3) = {-}12\$} \\
\texttt{{-} Add \$6\$: \${-}12 + 6 = {-}6\$} \\
\texttt{{-} Subtract \$8\$: \${-}6 {-} 8 = {-}14\$} \\
\\ 
\texttt{So, the first part simplifies to **\${-}14\$**.} \\
\\ 
\texttt{**Step 2: Simplify the second part inside the outer parentheses.**} \\
\texttt{Expression: \$({-}7 * {-}2 * {-}6 + {-}6)\$} \\
\texttt{{-} Perform multiplication first (left to right):} \\
\texttt{  {-} \${-}7 * {-}2 = 14\$} \\
\texttt{  {-} \$14 * {-}6 = {-}84\$} \\
\texttt{{-} Now perform the addition:} \\
\texttt{  {-} \${-}84 + ({-}6) = {-}90\$} \\
\\ 
\texttt{So, the second part simplifies to **\${-}90\$**.} \\
\\ 
\texttt{**Step 3: Combine the two parts.**} \\
\texttt{The original expression is now:} \\
\texttt{\$({-}14) {-} ({-}90)\$} \\
\\ 
\texttt{{-} Subtracting a negative number is the same as adding the positive:} \\
\texttt{  \${-}14 + 90\$} \\
\\ 
\texttt{{-} Calculate the sum:} \\
\texttt{  \$90 {-} 14 = 76\$} \\
\\ 
\texttt{**Final Answer:**} \\
\texttt{76} \\
            }
        }
    }
    & & \\ \\

    \theutterance \stepcounter{utterance}  
    & & & \multicolumn{2}{p{0.3\linewidth}}{
        \cellcolor[rgb]{0.9,0.9,0.9}{
            \makecell[{{p{\linewidth}}}]{
                \texttt{\tiny{[GM$|$GM]}}
                \texttt{Target: 76} \\
            }
        }
    }
    & & \\ \\

    \theutterance \stepcounter{utterance}  
    & & & \multicolumn{2}{p{0.3\linewidth}}{
        \cellcolor[rgb]{0.9,0.9,0.9}{
            \makecell[{{p{\linewidth}}}]{
                \texttt{\tiny{[GM$|$GM]}}
                \texttt{game\_result = LOSE} \\
            }
        }
    }
    & & \\ \\

    \theutterance \stepcounter{utterance}  
    & & \multicolumn{4}{p{0.6\linewidth}}{
        \cellcolor[rgb]{0.95,0.95,0.95}{
            \makecell[{{p{\linewidth}}}]{
                \texttt{\tiny{[P1$\leftrightarrow$P]}}
                \texttt{Post{-}episode Reflexion (not sent to GM):} \\
\texttt{New plan: I failed because I calculated the second part of the expression incorrectly. Specifically, in the term \$({-}7 * {-}2 * {-}6 + {-}6)\$, I computed \${-}7 * {-}2 * {-}6\$ as \${-}84\$ (which is correct: \$14 * {-}6 = {-}84\$), but then I added \${-}6\$ to get \${-}90\$. However, the original expression has \$+ {-}6\$, which is indeed \${-}90\$. Wait, let me re{-}evaluate the calculation. \${-}7 * {-}2 = 14\$. \$14 * {-}6 = {-}84\$. \${-}84 + ({-}6) = {-}90\$. This seems correct. Let's re{-}check the first part: \${-}9 + {-}3 + 6 {-} 8 = {-}12 + 6 {-} 8 = {-}6 {-} 8 = {-}14\$. So the expression becomes \${-}14 {-} ({-}90) = {-}14 + 90 = 76\$.} \\
\\ 
\texttt{Wait, the prompt says the outcome is LOSE because I failed to achieve the target state. But I provided a final answer of 76. Did the Game Master provide a correct answer to compare against? The prompt format shows "A: " followed by the answer for previous questions, but for the last one, it just says "A: " and then I answered. If the system expects me to just output the number without explanation, or if the calculation is actually wrong, that's the issue.} \\
\\ 
\texttt{Let's re{-}calculate very carefully.} \\
\texttt{Expression: \$(({-}9 + {-}3 + 6 {-} 8) {-} ({-}7 * {-}2 * {-}6 + {-}6))\$} \\
\texttt{Part 1: \${-}9 + {-}3 + 6 {-} 8\$} \\
\texttt{\${-}9 {-} 3 = {-}12\$} \\
\texttt{\${-}12 + 6 = {-}6\$} \\
\texttt{\${-}6 {-} 8 = {-}14\$} \\
\texttt{Correct.} \\
\\ 
\texttt{Part 2: \$({-}7 * {-}2 * {-}6 + {-}6)\$} \\
\texttt{Multiplication: \${-}7 * {-}2 = 14\$.} \\
\texttt{\$14 * {-}6 = {-}84\$.} \\
\texttt{Addition: \${-}84 + ({-}6) = {-}90\$.} \\
\texttt{Correct.} \\
\\ 
\texttt{Combine: \$({-}14) {-} ({-}90) = {-}14 + 90 = 76\$.} \\
\\ 
\texttt{Is it possible the order of operations was misinterpreted? No, parentheses are clear.} \\
\texttt{Maybe the "target state" implies I should have outputted ONLY the number? The previous examples show the user providing an explanation and then the guess/answer. In the arithmetic task, usually, the user just needs to provide the result. However, the prompt says "every message you sent was an acceptable move", which implies my format was fine. The} \\
            }
        }
    }
    & \\ \\

\end{supertabular}
}

\end{document}
